
\section{Where Can Anyone Get Bread?}

Mark 8:1

“Folkebladet,” August 01, 1917

There is much anxiety contained in this question. Many are those who, with anguished hearts and tear-filled or, still worse, tearless eyes, have asked in this way. And many are still those who ask this anxious question. Those who formed the question in the biblical passage from which the heading of this little piece is taken addressed it to Jesus, the Savior. And we know that with him help was to be found.

It is rather interesting to notice how Jesus regarded human bodily distress and need. His first miracle was to remove the embarrassment of a poor bridal couple. And we know that on several occasions he showed that he “was deeply moved with compassion” for the people because of their bodily distress. And as he was, so he still is. “Jesus Christ is the same yesterday and today, yes, forever!”

It is good to remember this also in the serious time in which we live, and in which the question, “Where can anyone get bread?” becomes very pressing for many. He still has the heart that “is deeply moved with compassion” for those who suffer bodily hunger, or who in some other way are oppressed by sorrow and anxiety. And he still possesses the power that enables him to feed the hungry and to comfort the sorrowful and unhappy.

Go, therefore, with the many anxious questions that stir within your soul to Jesus. With him help is to be found, peace and rest to be found. Say to God with childlike confidence:

“Care, O dear Father, you!
I will not be anxious,
nor with a troubled mind
ask about my future.
Care for me throughout my time,
care for me and mine!
God, almighty, gracious, gentle,
care for all who are yours!

Care when I go to rest
and my eye closes.
Care when I rise again
and hasten to my work!
Care thus for my calling and station,
hand and mouth and heart,
for the work that I can do,
for my joy and pain!”

If Jesus had lived bodily among us, he would surely have wished to teach us the same thing that through his Word he will teach us: to live a simple and frugal life. He feeds the hungry crowds there in the wilderness, it is true. But he feeds them with bread and fish—and nothing more. Thus, not so many dishes, so many kinds of food as are now so commonly used, but which are in no way necessary in order to eat one’s fill. And he would teach us to “gather the remnants so that nothing is wasted.” In our time there is a boundless amount that “goes to waste,” and we greatly need to learn from Jesus to take care of what remains and to learn contentment with little.

Besides this, he will teach us to “give thanks.” It is written that Jesus “took the seven loaves and gave thanks.” He therefore prayed at table. There are many who do not do this, including many of our Scandinavian people in this country. In this respect they do not resemble him whose name they nevertheless bear and wish to bear. More than once I have found myself deeply grieved that even such young people as have come from Christian homes, where prayer at table is regularly offered by both parents and children, as soon as they have been out among strangers, cease to do as Jesus did: “give thanks.”

Why? Ah, because “the others” do not do it; the Americans do not do it; it is not refined enough. And when these same young people come home to father and mother, where prayer is still offered at the table, they are ashamed to fold their hands and join their praying parents in the table prayer. I would like to be permitted to say to young and old: let us do as Jesus did—give thanks for our daily bread and ask the Lord to bless it for us.

Besides this, Jesus will teach us in this Word that he desires it, wills it, that his children, his people, are to help God feed others, just as the disciples helped Jesus give food to the people.

This still applies. And it applies not only in the question of food for the body. It applies especially in the question of that food by which the soul may be filled and live.

Jesus Christ wills that we should always carry the picture of our poor brother’s need and distress in our soul, so that we are driven to share with him:

“Up, brothers, hasten—the life is short!
Much should be done, but so little is done.
Many are those who wait—those who give are but few;
enough there is for the great, but so little for the small.
Wipe the tear away, brothers, wipe tears away!
Life is short from our cradle to grave.
The Lord, he waits; we are to follow in his footsteps;
God will bless everyone who blesses his brother!”

And here I must then mention our home mission. Those who have been appointed to care for it often ask in deep anxiety: “Where can anyone get bread?” when our workers do not receive their support, when the mission treasury is empty, when the debt grows, and when little or nothing comes in by way of gifts for this work. The Mission Board also tries to take this anxious question to Jesus. But when he answers now as then, “How many loaves do you have?” he does not say this only to us who have been appointed to direct the work of our home mission, but he says it to all of you—and also to you, my brother and sister, who read this.

And what shall we answer him? I know what we can answer him if only we are willing: We have enough bread for the need, provided it is distributed with Jesus’ blessing upon it, if only we are willing to distribute it to those who are in need.

It is written that Jesus “gave them to his disciples so that they should set them before the people.” They were not merely to have the bread for themselves—they were to distribute it to the people.

Yes, my brother and sister, do you understand what Jesus would have you do? You have bread—food and drink and God’s saving Word—in abundance. Distribute it to your poor, needy, hungry brothers! There is blessing in it.

E. B.